% Case Study Figure — paste into BHI 2026 paper.tex
% Required packages (add to preamble if not already present):
%   \usepackage{tcolorbox}
%   \tcbuselibrary{skins,breakable}
%   \usepackage{tabularx}
%   \usepackage{makecell}

\begin{figure*}[t]
\centering
\begin{tcolorbox}[
  enhanced,
  title={\textbf{Case Study: End-to-End Pipeline Output} \hfill
         \small\textit{Patient sub-023\_S\_6334} $\cdot$
         78-year-old male $\cdot$ 14 years education $\cdot$
         \textbf{Clinical label: MCI}},
  colframe=black!70,
  colback=white,
  coltitle=white,
  colbacktitle=black!70,
  fonttitle=\small,
  top=4pt, bottom=4pt, left=4pt, right=4pt,
  boxrule=0.6pt,
  arc=3pt
]

%% ── Row 1: Classification  +  KG Retrieval ─────────────────────────────
\begin{minipage}[t]{0.48\textwidth}
\begin{tcolorbox}[
  enhanced, colframe=gray!50, colback=gray!5,
  title={\small\textbf{(1) PCAG-ComBat Classification}},
  coltitle=black, fonttitle=\small\bfseries,
  top=2pt, bottom=2pt, left=3pt, right=3pt, boxrule=0.4pt, arc=2pt]
\small
\begin{tabular}{@{}lcc@{}}
\toprule
Task & Prob. & Pred.\\
\midrule
NC vs.\ AD  & 0.318 & NC\\
NC vs.\ MCI & 0.623 & \textbf{MCI}\,$\checkmark$\\
MCI vs.\ AD & 0.115 & MCI\\
\bottomrule
\end{tabular}\\[2pt]
\textbf{OVO Verdict:} \colorbox{yellow!40}{\textbf{MCI}} (medium confidence)
\end{tcolorbox}
\end{minipage}
\hfill
\begin{minipage}[t]{0.48\textwidth}
\begin{tcolorbox}[
  enhanced, colframe=gray!50, colback=gray!5,
  title={\small\textbf{(2) Knowledge Graph + RAG Retrieval}},
  coltitle=black, fonttitle=\small\bfseries,
  top=2pt, bottom=2pt, left=3pt, right=3pt, boxrule=0.4pt, arc=2pt]
\small
\textbf{Similar patients (Neo4j KG):}\\
\quad sub\_0117 (MCI, sim=0.818),\; 011\_S\_7048 (NC, sim=0.805)\\[2pt]
\textbf{Retrieved literature (MMR, $k{=}5$):}\\
\quad Vemuri et al.\ (2014) — Cognitive Reserve \& MCI\\
\quad Zhao et al.\ (2022) — fMRI GCN for dementia prediction\\
\quad Jack et al.\ (2010) — AD biomarker cascade model
\end{tcolorbox}
\end{minipage}

\vspace{4pt}

%% ── Row 2: Neural Evidence (sMRI + fMRI) ────────────────────────────────
\begin{tcolorbox}[
  enhanced, colframe=gray!50, colback=gray!5,
  title={\small\textbf{(3) Multimodal Neural Evidence}},
  coltitle=black, fonttitle=\small\bfseries,
  top=2pt, bottom=2pt, left=3pt, right=3pt, boxrule=0.4pt, arc=2pt]
\small
\begin{minipage}[t]{0.47\textwidth}
\textbf{sMRI — BrainIAC ViT Attention Rollout:}\\[1pt]
\begin{tabular}{@{}lc@{}}
ROI & Attention\\
\hline
Cingulum\_Mid\_L & 0.288\\
Cingulum\_Post\_R & 0.155\\
Precuneus\_R & 0.142\\
Cingulum\_Post\_L & 0.144\\
Caudate\_L & 0.070\\
\end{tabular}
\end{minipage}
\hfill\vline\hfill
\begin{minipage}[t]{0.47\textwidth}
\textbf{fMRI — GAT Network-Level Attention:}\\[1pt]
\begin{tabular}{@{}lc@{}}
ROI & Salience\\
\hline
Cingulum\_Mid\_R & 1.000\\
Temporal\_Sup\_R & 0.973\\
Precuneus\_R & 0.953\\
Cingulum\_Mid\_L & 0.961\\
Rolandic\_Oper\_L & 0.978\\
\end{tabular}
\end{minipage}
\end{tcolorbox}

\vspace{4pt}

%% ── Row 3: Generated Report Excerpt ─────────────────────────────────────
\begin{tcolorbox}[
  enhanced, colframe=gray!50, colback=blue!3,
  title={\small\textbf{(4) Gemma3-12B Generated Report (excerpt, translated)}},
  coltitle=black, fonttitle=\small\bfseries,
  top=2pt, bottom=2pt, left=3pt, right=3pt, boxrule=0.4pt, arc=2pt]
\small\itshape
``BrainIAC attention concentrates on bilateral Cingulum (mid/post, attn$\,{=}\,$0.29--0.14)
and Precuneus (0.14), structural regions implicated in early MCI and default-mode network
degeneration. GAT functional analysis reveals concordant disruption in Cingulum and
superior temporal cortex (salience ${\geq}$0.95), consistent with known salience-network
and language-association-area vulnerability in MCI [Vemuri et al., 2014].
Structural and functional evidence are mutually reinforcing, supporting the
\textbf{MCI} classification.
Given the patient's educational attainment (14 years), preserved cognitive reserve may
partially compensate for early network disruption.
\textbf{Recommendation:} Longitudinal follow-up with cognitive assessment every 6--12 months;
monitor Cingulum atrophy rate as a progression biomarker.''
\end{tcolorbox}

\end{tcolorbox}

\caption{End-to-end pipeline case study. For a correctly classified MCI patient, the system
integrates PCAG-ComBat predictions (Step~1), KG-retrieved similar patients and
RAG-retrieved literature (Step~2), multimodal neural evidence from BrainIAC ViT
attention rollout and GAT network-level attention (Step~3), and a Gemma3-12B generated
diagnostic report grounded in all upstream evidence (Step~4).
Concordant Cingulum and Precuneus involvement across both modalities supports the MCI staging.}
\label{fig:casestudy}
\end{figure*}
